\documentclass[10pt,a4paper]{amsart}
\usepackage[margin=19mm]{geometry}
\usepackage{amsmath,amssymb,mathtools}
\usepackage[scr=rsfs,cal=euler]{mathalfa}
\usepackage{xfrac}
\usepackage[unicode,hidelinks,bookmarks=false]{hyperref}
\newcommand{\ie}{i.e.\ }
\newcommand{\cf}{cf.\ }
\newcommand{\resp}{respectively\ }
\newcommand{\Subsection}{\subsection}
\providecommand{\nobreakdash}{\nobreak\hbox{-}}
\setlength{\emergencystretch}{2em}
\multlinegap0pt
\usepackage{amssymb}
\usepackage{tikz}
\usepackage{bm}
\newtheorem{prop}{Proposition}
\newtheorem{thm}{Theorem}
\newtheorem{lem}{Lemma}
\def\A{\mathbb A}
\def\C{\mathbb C}
\def\EE{\mathcal E}
\def\GL{{\rm GL}}
\def\Q{\mathbb Q}
\def\R{\mathbb R}
\def\Res{{\rm Res}}
\def\g{\mathfrak g}
\def\h{\mathfrak h}
\def\m{\mathfrak m}
\def\vv{{\sf v}} %harish-chandra or infinity type
\def\ww{\mathsf w}
\title{On the residual Eisenstein cohomology of unitary groups}
\author{Harald Grobner}
\date{Source: arXiv:2608.15947v2 (CC BY 4.0)}
\begin{document}
\maketitle

\noindent\emph{Source and licence.} On the residual Eisenstein cohomology of unitary groups. Version arXiv:2608.15947v2 (CC BY 4.0), distributed by arXiv under the Creative Commons Attribution 4.0 International licence, http://creativecommons.org/licenses/by/4.0/, as recorded in the arXiv licence metadata for this exact version and shown on the abstract page https://arxiv.org/abs/2608.15947v2. Full licence notice: \texttt{licenses/arxiv-2608.15947-CC-BY-4.0.txt}.\par\smallskip

\noindent\textit{Editorial scope.} Counted unit: the article's thm prop:nonempty-P2-example (source0.tex lines 1028--1044). The article's complete original proof of it is delivered with it (source0.tex lines 1045--1210, one \texttt{proof} environment of 166 lines). No mathematics is written, repaired or re-derived: the statement and the proof are the article's own lines, in the article's own notation, and no retained line is substituted or re-flowed.  Retained uncounted as the source-local context the proof needs: lines 258--258; lines 363--366; lines 386--393; lines 418--418; lines 578--598; lines 647--666; lines 676--676; lines 719--719; lines 904--914; lines 980--1022.  Nothing was omitted from any retained span, so the package carries no omission record.  The retained lines cite 8 works, whose entries are transcribed verbatim from the article's own bibliography source in the e-print and delivered with short editorial labels recorded in the item's reference mapping.  No retained line contains a figure environment, an image-inclusion command or drawing code, so no image is delivered and the package declares no asset dependency.  Editorial formatting only, in addition: an AMS \texttt{amsart} preamble that restates the article's own declarations and macros used by the retained lines, namely the environments \texttt{prop}, \texttt{thm}, \texttt{lem} on separate counters, together with the standard packages those lines need.  The retained environments print in this document as Proposition 1, Theorem 1, Lemma 1 in order of appearance, whereas the article prints the same items with the numbering of its own full document.  The complete native source of the article as distributed in the arXiv e-print for this exact version is bundled unchanged as \texttt{sources/207798/source0.tex}.
\subsection{Number fields}\label{sect:nf}

\begin{prop}\label{prop:qminLk}
Let $\pi_\infty=\tau_\infty\hat\otimes\sigma_\infty$ be a cohomological irreducible tempered representation of $L_{k,\infty}=(\Res_{E/\Q}\GL_n)(\R)\times (\Res_{F/\Q})(U(V_{n-2k}))(\R)$. Then
$$q_{min}(\pi_\infty)=\sum_{\vv\in S^{\sf in}_\R}(p_\vv-k)(q_\vv-k) + |S^{\sf in}_\R| \cdot\tfrac{k(k-1)}{2}+ |S^{\sf sp}_\R| \cdot\left(\big\lfloor\tfrac{n-2k}{2}\big\rfloor\big\lceil\tfrac{n-2k}{2}\big\rceil+2\big\lfloor\tfrac{k}{2}\big\rfloor\big\lceil\tfrac{k}{2}\big\rceil\right) + |S_\C| \cdot\tfrac{n^2-4kn-n+6k^2}{2}$$
\end{prop}

\begin{thm}\label{thm:qtcusp}
Let $\pi$ be a cohomological, square-integrable automorphic representation of $G_n(\A_F)$ and write $\Pi=\Pi_1\boxplus...\boxplus\Pi_l$ for its base change to $\GL_n(\A_E)$. Then the following assertions are equivalent:
\begin{enumerate}
\item The summands $\Pi_i$ are cuspidal for all $1\leq i\leq l$.
\item $\Pi$ is tempered at all archimedean places and quasi-tempered at all non-archimedean places $\ww\in\Sigma$. Here, the latter means that for each non-archimedean $\ww\in\Sigma$, there is a partition $n=m_1+...+m_u$, discrete series representations $\delta_{i}$ of $\GL_{m_i}(E_\ww)$, and real numbers $a_{i}$ with $|a_{i}|<\tfrac12$, such that $\Pi_\ww\cong\delta_{1} |.|_\ww^{a_{1}}\times...\times\delta_{u} |.|_\ww^{a_{u}}$. 
\end{enumerate}
If one of the above equivalent assertions holds, then $\pi$ is necessarily cuspidal, $\pi_\infty$ is tempered and the isobaric summands $\Pi_i$ are all distinct, $\Pi_i\ncong\Pi_j$, for $i\neq j$. Moreover, for each $1\leq i\leq l$, the partial Asai $L$-function $L^T(s,\Pi_i,{\rm As}^{(-1)^{n-1}})$ has a pole at $s=1$ for a sufficiently large, but finite set of places $\Sigma_\infty\subseteq T\subset \Sigma$.
\end{thm}

\subsection{Parabolic and cuspidal supports}\label{sect:pc}

\begin{lem}\label{lem:Kostant}
Let $P_k$ be the standard maximal parabolic $F$-subgroup of $G$ corresponding to $\alpha_k\in\Delta^0_F$, $1\leq k\leq r$. Let $\vv\in S_\R$. Then, $W_\vv=W(\g_{\vv,\C},\h_{\vv,\C})\cong \mathfrak S_n$ may be identified with the symmetric group of permutations of the set $\{1,...,n\}$ and 
\begin{equation}\label{eq:Kostant}
W_\vv^{P_k}=
\left\{
w \in \mathfrak S_n \;\middle|\;
\begin{aligned}
& w^{-1}(1) < \cdots < w^{-1}(k), \\
& w^{-1}(k+1) < \cdots < w^{-1}(n-k), \\
& w^{-1}(n-k+1) < \cdots < w^{-1}(n)
\end{aligned}
\right\}.
\end{equation}
Hence, $W_\vv^{P_k}$ is parametrized as a subset of $\mathfrak S_n$ by pairs $(I,J)$ of disjoint subsets
$I:=\{i_1, \cdots, i_k\}\subseteq\{1,\dots,n\}$ and $J:=\{j_1,...,j_{k}\}\subseteq\{1,\dots,n\}$, with $i_1<...<i_k$ and $j_1<...<j_k$, and complement
$\{1,\dots,n\}\setminus (I\cup J)=:\{r_1, ..., r_{n-2k}\}$, with $r_1< ...< r_{n-2k}$ via
$$
w_{I,J}(i_\ell)=\ell,\qquad w_{I,J}(j_\ell)=n-k+\ell, \qquad w_{I,J}(r_\ell)=k+\ell.
$$
If $\vv\in S_\C$, then $W_\vv^{P_k}$ is given by the product of two identical, but independent copies of sets of type \eqref{eq:Kostant}.
\end{lem}

\begin{equation}\label{sw}
\boxed{s_{w}
=
\begin{cases}
\displaystyle
\frac{1}{2k}
\left(\sum_{j_\vv\in J_\vv}(\mu_{{\vv},j_\vv}-j_\vv)
-
\sum_{i_\vv\in I_\vv}(\mu_{{\vv},i_\vv}-i_\vv)\right),
&\forall \vv\in S_\R,\\[16pt]
\displaystyle
\frac{1}{4k}
\left(\sum_{j_{\iota_\vv}\in J_{\iota_\vv}}(\mu_{{\iota_\vv},j_{\iota_\vv}}-j_{\iota_\vv})+\sum_{j_{\bar\iota_\vv}\in J_{\bar\iota_\vv}}(\mu_{{\bar\iota_\vv},j_{\bar\iota_\vv}}-j_{\bar\iota_\vv})
-
\left(\sum_{i_{\iota_\vv}\in I_{\iota_\vv}}(\mu_{{\iota_\vv},i_{\iota_\vv}}-i_{\iota_\vv})
+
\sum_{i_{\bar\iota_\vv}\in I_{\bar\iota_\vv}}(\mu_{{\bar\iota_\vv},i_{\bar\iota_\vv}}-i_{\bar\iota_\vv})\right)\right),
& \forall\vv\in S_\C.
\end{cases}}
\end{equation}\normalsize

\subsubsection{}\label{sect:extex}

\subsubsection{Tables of representatives}\label{sect:tablesw}

\begin{thm}\label{thm:main}
Let $E/F$ be a quadratic extension of arbitrary number fields and let $G=G_n$ be the unitary group over $F$, attached to a non-dengenerate hermitian form on an $n$-dimensional $E$-vector space of $F$-rank $r\geq 1$. Let $P_k$, $1\leq k\leq r$, be a standard maximal parabolic $F$-subgroup of $G$ and let $\varphi_{P_k} $ be an associate class of cuspidal automorphic representations of $L_k(\A_F)\cong \GL_k(\A_E)\times G_{n-2k}(\A_F)$, cf.\ Sect.\ \ref{sect:pc}, which is represented by a cuspidal automorphic representation $\pi_{\lambda_0}=\pi\otimes e^{\langle \lambda_0, H_P(\cdot)\rangle}$. We write, as we may, $\pi=\tau\hat\otimes\sigma$ for a unitary cuspidal automorphic representation $\tau$ of $\GL_k(\A_E)$ and a unitary cohomological cuspidal automorphic representation $\sigma$ of $G_{n-2k}(\A_F)$ and we let  $w\in W^{P_k}$ be the corresponding unique Kostant representative such that $\pi_\infty$ is $(\m_{k,\infty},K_{M^\circ_{k,\infty}})$-cohomological with respect to $\EE_{\mu_w}|_{M^\circ_{k,\infty}}$ and such that $\lambda_0=s_w\cdot \gamma_k$ with $s_w$ being given by \eqref{sw}. Furthermore, we abbreviate
$$q(\pi_{\lambda_0}):=q_{min}(\pi_\infty) + [F:\Q] k(2n-3k) -\ell(w),$$
where $q_{min}(\pi_\infty)$ is given by Prop.\ \ref{prop:qminLk}.
We assume that $\Pi=BC(\sigma)$ satisfies the equivalent conditions of Thm.\ \ref{thm:qtcusp}, i.e., $\Pi=\Pi_1\boxplus...\boxplus\Pi_l$ for cuspidal automorphic representations $\Pi_i$, $1\leq i\leq l$ and that either
\begin{enumerate}
\item $s_w=\tfrac12$ and $\tau$ is conjugate self-dual, $\epsilon$-distinguished and $L(\tfrac12,\tau\times \Pi^\vee)\neq 0$; or
\item $s_w=1$ and $\tau=\Pi_i$ for some $1\leq i\leq l$.
\end{enumerate}
Then, $r^{q(\pi_{\lambda_0})}_{\varphi_{P_k}}$ is injective and its image, which consists of residual Eisenstein cohomology classes, is non-zero.
\end{thm}

We resume the discussion of our extended example $G=U(V_6,h)/\Q(\sqrt[3]{2})$ from Sect.\ \ref{sect:extex} in the light of Thm.\ \ref{thm:main}. Our aim is to give concrete examples of non-zero residual Eisenstein cohomolgy classes for $G(\A_F)$, by making all ingredients of our general theorem Thm.\ \ref{thm:main} explicit (and showing their existence). Firstly, we note that
$$
\dim_\mathbb R(G_\infty/K_G)
=
\dim_\mathbb R U(3,3)/ (U(3)\times U(3))
+
\dim_\mathbb R \GL_6(\mathbb C)/U(6)
=
2\cdot 3\cdot 3+6^2
=
54.
$$
and we summarize the calculated values of several input-data in the following table:
\begin{table}[ht]
\centering
\small
\renewcommand{\arraystretch}{1.2}
\begin{tabular}{c|c|c|c|c|c|c}
\hline
\(P_k\)
&
\(s_w\)
&
\(\begin{array}{c} q_{min}(\pi_\infty) \end{array}\)
&
\(\begin{array}{c} q(\pi_{\lambda_0}) \end{array}\)
&
\(\begin{array}{c} \dim_\mathbb R X_{L_k} \end{array}\)
&
\(\begin{array}{c} \dim_\mathbb R N_{k,\infty} \end{array}\)
&
\(\begin{array}{c} \dim_\mathbb R \partial_{P_k}(\overline X_G) \end{array}\)
\\
\hline
\(P_1\) & \(\frac12\) & \(10\) & \(21,\ 22\)             & \(26\) & \(27\) & \(53\)\\
\(P_2\) & \(\frac12\) & \(5\)  & \(17,\ 18,\ 19,\ 20\)   & \(17\) & \(36\) & \(53\)\\
\(P_2\) & \(1\)       & \(5\)  & \(14,\ 15,\ 16,\ 17\)   & \(17\) & \(36\) & \(53\)\\
\(P_3\) & \(\frac12\) & \(9\)  & \(18\)                  & \(26\) & \(27\) & \(53\)\\
\hline
\end{tabular}
\caption{Cohomological degrees and dimensions of related Borel--Serre boundary data}
\label{tab:example-degrees-boundary}
\end{table}

\begin{thm}\label{prop:nonempty-P2-example}
Let $w\in W^{P_2}$ be the Kostant representative $w=
        \bigl(
        w_{\vv_\R},w_{\vv_\C}
        \bigr)
        =
        \bigl(
        w_0,(w_0,w_0)
        \bigr)$, with $w_0=(5,3,1,6,4,2)$
occurring in the column of \(s_w=1\) in the second table of Sect.\ \ref{sect:tablesw}. Then there
exists a cuspidal automorphic representation $\pi=\tau\widehat\otimes\sigma$ of $L_2(\mathbb A_F)
        \cong
        \GL_2(\mathbb A_E)\times U(V_2,h)(\mathbb A_F)$
satisfying the hypotheses of Theorem~\ref{thm:main} at $\lambda_0=\gamma_2.$
Consequently, degree $ q(\pi_{\lambda_0})
        = 14$ in the third row of Table \ref{tab:example-degrees-boundary} carries non-zero residual Eisenstein cohomology classes. None of them lies in inner cohomology.
\end{thm}

\begin{proof}
By our recipe from Lem.\ \ref{lem:Kostant},  the pair $(I,J)$, given by
$$
        I=\{3,6\},
        \qquad
        J=\{1,4\}.
$$
determines $w_0$. We let $\sigma_\infty$ be a fixed cohomological tempered representation of
$$
        U(V_2,h)_\infty
        \cong
        U(1,1)\times \GL_2(\mathbb C)
$$
determined by the complement $\{2,5\}$ of this Kostant datum. More precisely, this means that, since $\mu=0$ in our example, we choose a discrete series representation of $U(1,1)$, which is cohomological with respect to the irreducible finite dimensional representation of highest weight $(1,-1)$ (for the very moment it is not relevant wether it is holomorphic or anti-holomorphic, which will become important only later) and form its tensor product with the (unique) tempered representation of $\GL_2(\C)$, which is cohomological with respect to the irreducible finite dimensional representation of highest weight $((1,-1),(1,-1))$. As above, we write $BC(\sigma_\infty)$ for its local archimedean base-change. Explicitly, it factors as the product 
\begin{equation}\label{BCok}
BC(\sigma_\infty)\cong \bigotimes_{\ww\in\Sigma_\infty}
        \left(z^{3/2}\bar z^{-3/2}\right)
        \times
        \left(z^{-3/2}\bar z^{3/2}\right)
\end{equation}
of three isomorphic fully induced representations of three copies of $\GL_2(\mathbb C)$. It follows from the particular choice of \(w\), that the representation 
$$
        \pi_\infty
        :=
        \tau_\infty\widehat\otimes\sigma_\infty
        \quad\text{with}\quad
        \tau_\infty:= BC(\sigma_\infty).
$$
is $(\m_{2,\infty}, K^\circ_{M_{2,\infty}})$-cohomological with respect to $\EE_{\mu_w}|_{M^\circ_{2,\infty}}$.\\\\ 
We now globalize $\tau_\infty$ and $\sigma_\infty$ to cuspidal automorphic representations $\tau$, respectively $\sigma$, which satisfy all the conditions of Thm.\ \ref{thm:main} for $\lambda_0=\gamma_2$. Once this is achieved, we obtain the desired conclusion of Thm.\ \ref{prop:nonempty-P2-example} by applying the main result Thm.\ \ref{thm:main} to the associate class $\varphi_{P_2}$ defined by $\pi_{\gamma_2}:=(\tau\hat\otimes\sigma)\otimes e^{\langle \gamma_2, H_{P_2}(\cdot)\rangle}$ and inspecting Table \ref{tab:example-degrees-boundary} for the explicit numerology concerning, e.g., $q(\pi_{\lambda_0})$.\\\\
To this end, let $L_1=\mathbb Q(\sqrt{-2})$. Denote by $c_1$ the only non-trivial automorphism of $L_1$ (i.e., complex conjugation). Let $\varepsilon_{L_1/\mathbb Q}: \mathbb Q^*\backslash\mathbb A_\mathbb Q^*\rightarrow\C^*$ be the quadratic Hecke character of attached to
\(L_1/\mathbb Q\) by class field theory and let $\eta_{L_1/\mathbb Q}: L_1^*\backslash \A_{L_1}^*\rightarrow\C^*$ by a ($c_1$-)conjugate self-dual extension of $\varepsilon_{L_1/\mathbb Q}$ as in Sect.\ \ref{sect:nf}. More concretely, one may pin down $\eta_{L_1/\mathbb Q}$ using an elliptic curve over
\(\mathbb Q\) with complex multiplication by \(\mathcal O_{L_1}\) and letting \(\psi\) be the algebraic Hecke character of \(L_1\) attached to such a CM elliptic curve by Deuring's construction, see, e.g., \cite[Chp.\ II, \S\S 9--10]{Sil94}. Then put $\eta_{L_1/\mathbb Q}
        :=
        \psi\,\|\cdot\|_{\mathbb A_{L_1}}^{-1/2}$. It follows that one has
$$
        \eta_{L_1/\mathbb Q,\infty}(z)
        =
        z^{1/2}\bar z^{-1/2}.
$$
Now define $\chi_{1,0}
        :=
        \eta_{L_1/\mathbb Q}^{\,3}.$ Then,
$$
        \chi_{1,0,\infty}(z)
        =
        z^{3/2}\bar z^{-3/2},
$$
while still
$$
        \chi_{1,0}
        \big|_{\mathbb A_\mathbb Q^\times}
        =
        \varepsilon_{L_1/\mathbb Q}^{\,3}
        =
        \varepsilon_{L_1/\mathbb Q},  \quad\quad {\rm and} \quad\quad   {}^{\,c_1} \chi_{1,0}=\chi_{1,0}^{-1}.
$$
We next modify \(\chi_{1,0}\) at one non-archimedean place, in order to obtain local regularity. Since \(L_1\) is linearly disjoint from the normal closure of \(E/\mathbb Q\), Chebotarev's theorem gives a prime number \(p\geq 3\), away from the finite conductor of \(\chi_{1,0}\), such that \(p\) splits completely in \(E\) and is inert in \(L_1\). Let $w_p$ be the unique place of \(L_1\) above \(p\) and choose a character
$$
        \beta:L_{1,w_p}^{*}\longrightarrow\mathbb C^{*}
$$
of odd order $b_\beta$ satisfying
$$
        \beta|_{\mathbb Q_p^*}=1 \quad\quad {\rm and}\quad\quad      \bigl(\chi_{1,0,w_p}\beta\bigr)^2\neq 1.
$$
Such a character exists because the unramified quotient $ L_{1,w_p}^{*}/\mathbb Q_p^*$ has arbitrarily ramified characters of \(p\)-power order. Since \(b_\beta\) is odd, choose
\(r\in\mathbb Z\) with $2r\equiv 1\pmod{b_\beta}$ and put
$$
        \alpha_{w_p}:=\beta^r.
$$
Then,
$$
        \alpha_{w_p}^2=\beta,
        \qquad
        \alpha_{w_p}|_{\mathbb Q_p^*}=1.
$$
By the Grunwald--Wang theorem (recall that $b_\beta$ and $p$ are odd), there exists hence a finite order Hecke character
$$
        \xi_0:L_1^*\backslash \mathbb A_{L_1}^{*}
        \longrightarrow \mathbb C^*
$$
such that
$$
        (\xi_0)_{w_p}=\alpha_{w_p} .
$$
See Artin--Tate, \cite{ArtinTate}, Chp.~X, Thm.~5, or the formulation in \cite{BCGNT}, Lem.\ 5.3.1.  Define
$$
        \xi:=\xi_0({}^{\,c_1}\xi_0)^{-1}.
$$
Then, $\xi$ is ($c_1$-)conjugate self-dual, i.e., ${}^{\,c_1}\xi=\xi^{-1}$, by construction and $\xi|_{\mathbb A_\mathbb Q^*}=1$, as \(c_1\) acts trivially on \(\mathbb A_\mathbb Q^*\). Moreover, at \(w_p\) we have $\xi_{w_p}=\alpha_{w_p}\,({}^{\,c_1}\alpha_{w_p})^{-1},$
which, since \(\alpha_{w_p}\) is trivial on \(\mathbb Q_p^*\) and therefore ${}^{\,c_1}\alpha_{w_p}=\alpha_{w_p}^{-1}$ (recall that $p$ is inert in $L_1$, whence $c_1$ is the non-trivial automorphism of $L_{1,w_p}/\Q_p$), simplifies to $\xi_{w_p} = \alpha_{w_p}^2 = \beta.$ Set
$$
        \chi_1:=\chi_{1,0}\cdot \xi.
$$
Then, $\chi_1$ is a ($c_1$-)conjugate self-dual Hecke character $L^*_1\backslash \A^*_{L_1}\rightarrow\C^*$ which satisfies
$$
        \chi_{1,\infty}(z)
        =
        z^{3/2}\bar z^{-3/2}, \quad\quad
        \chi_1|_{\mathbb A_\mathbb Q^*}
        =
        \varepsilon_{L_1/\mathbb Q},\quad\quad {\rm and} \quad\quad
        \chi_{1,w_p}^2\neq 1.
$$
Now put $L:=EL_1$ and let $\theta$ be the unique non-trivial element of \(\operatorname{Gal}(L/E)\). Thus, \(\theta\) acts trivially on \(E\) and as complex conjugation on \(L_1\). Let \(\widetilde c\) be the automorphism of \(L\) which acts as the non-trivial element \(c\in\operatorname{Gal}(E/F)\) on \(E\) and trivially on \(L_1\). Then \(\theta\) and \(\widetilde c\) commute. Define a Hecke character of \(L\) by
$$
        \chi:=\chi_1\circ N_{L/L_1}.
$$
Then, the following holds:
$$\boxed{{}^{\theta}\chi=\chi^{-1}\neq\chi \quad\quad {\rm and} \quad\quad {}^{\widetilde c}\chi=\chi}$$ 
Indeed, $N_{L/L_1}(\theta x)
        =
        c_1\bigl(N_{L/L_1}(x)\bigr)$ and since \({}^{c_1}\chi_1=\chi_1^{-1}\), it follows that ${}^\theta\chi=\chi^{-1}$, which shows the first claimed equation of the box. In order to see the inequality, mentioned there, we choose a place \(\ww_p\) of \(E\) above \(p\) and
a place \(u_p\) of \(L\) above \(\ww_p\). Since \(p\) splits completely in \(E\), one has
$$
        E_{\ww_p}\cong \mathbb Q_p,
        \qquad
        L_{u_p} \cong L_{1,w_p},
$$
whence under this identification the local component of
\(\chi=\chi_1\circ N_{L/L_1}\) is $\chi_{u_p}=\chi_{1,w_p}.$ As by construction, $\chi_{1,w_p}^2\neq 1$, 
$$
        {}^\theta\chi_{u_p}=\chi_{u_p}^{-1}\neq \chi_{u_p}.
$$
In particular, globally, $\chi^{-1}\neq \chi$. In order to prove the last claim of the above box, ${}^{\widetilde c}\chi=\chi$, just observe that \(\widetilde c\) acts trivially on \(L_1\), and that the norm \(N_{L/L_1}\) is invariant under automorphisms over \(L_1\).\\\\ With this Hecke character $\chi$ at hand, let
$$
        \Pi:=\operatorname{AI}_{L/E}(\chi)
$$
be the {\it automorphic induction} of $\chi$ to \(\GL_2(\mathbb A_E)\) as it has been established in the form needed here in Chp.\ 3, Thm.\ 6.2 of \cite{arthur-clozel} (see also \cite{henniart12}, Thm.\ 2). Since ${}^\theta\chi\neq\chi$, as we have just observed, \(\Pi\) is a cuspidal automorphic representation of \(\GL_2(\mathbb A_E)\) by  \cite{arthur-clozel}, Lem.\ 6.4 in Chp.\ 3. We claim that $\Pi$ is conjugate self-dual (with respect to $c$): Indeed, on the one hand, as ${}^\theta\chi=\chi^{-1}$,
$$ \Pi^\vee
        \cong
        \operatorname{AI}_{L/E}(\chi^{-1})
        =
        \operatorname{AI}_{L/E}({}^\theta\chi)
        \cong
        \Pi$$
but also, since  \( {}^{\widetilde c}\chi=\chi\), 
$$
        {}^c\Pi
        \cong
        \operatorname{AI}_{L/E}({}^{\widetilde c}\chi)
        =
        \operatorname{AI}_{L/E}(\chi)
      =
        \Pi
$$
whence in summary
$$
        {}^c\Pi\cong\Pi^\vee
$$
as claimed. At the archimedean places, the local extension \(L/E\) is split, and
automorphic induction gives the sum of the two characters obtained from
\(\chi_1\) and \({}^c\chi_1\). Therefore,
$$
        \Pi_\infty\cong \bigotimes_{\ww\in\Sigma_\infty}
        \left(z^{3/2}\bar z^{-3/2}\right)
        \times
        \left(z^{-3/2}\bar z^{3/2}\right),
$$
whence by \eqref{BCok} $\Pi_\infty\cong BC(\sigma_\infty)$ and we set as the first of our two desired globalizations
$$
\tau:=\Pi.       
$$
In order for this choice of a cuspidal automorphic representation $\tau$ of \(\GL_2(\mathbb A_E)\) to be suited to be the \(\GL_2(\mathbb A_E)\) -factor of a unitary cuspidal automorphic rerpesentation $\pi=\tau\hat\otimes\sigma$ of $L_2(\A_F)$, which matches all the conditions of Thm.\ \ref{thm:main} in the case of $\lambda_0=\gamma_2$, it is therefore now necessary and sufficient to check that $\tau$ descends to a cuspidal automorphic representation of $U(V_2,h)(\A_F)$ (and then to let $\sigma$ be this descent): In fact, if $\tau=BC(\sigma')$ with $\sigma'$ a cuspidal automorphic representation of $U(V_2,h)(\A_F)$, then the archimedean component $\sigma'_\infty$ is a tempered cohomological representation of $U(V_2,h)_\infty\cong U(1,1)\times \GL_2(\C)$ of the right highest weight determined by the fixed Kostant representative $w=(w_0,(w_0,w_0))$, as so was its local base change $\tau_\infty$. At the cost of changing the cohomological discrete series representation at the $U(1,1)$-factor of our previously fixed representation $\sigma_\infty$ from the beginning of this proof, we may assume that $\sigma_\infty=\sigma'_\infty$ and hence set $\sigma:=\sigma'$ globally. \\\\ 
We are hence left to show that such a cuspidal automorphic representation $\sigma'$ exists. For that, we first observe that by \cite{mok}, Cor.\ 2.5.9, $L^T(s,\Pi,\operatorname{As}^{-})$ has a pole at $s=1$, for a sufficiently big finite set of places $T$ of $F$. Therefore, $\Pi$ descends to $U(V_2,h)(\A_F)$ by \cite{mok}, Thm.\ 2.5.4, (see also his Rem.\ 2.5.5. Alternatively, the reader may have a look at \cite{zou}, Thm.\ 2.1 and Thm.\ 2.6 {\it ibidem}) i.e., there is an irreducible square-integrable automorphic representation $\sigma'$ of $U(V_2,h)(\A_F)$, whose stable base change is $\Pi=BC(\sigma')$. Such a representation $\sigma'$ is cuspidal as desired: Indeed, as the local Weil representation $\operatorname{Ind}_{W_{L_{u_p}}}^{W_{E_{\ww_p}}}\chi_{u_p}$ is irreducible, the local component $\Pi_{\ww_p}$, being the local automorphic induction, is supercuspidal. Consequently, since at a split non-archimedean place \(\vv_p\) of \(F\) below \(\ww_p\), we have $U(V_2,h)(F_{\vv_p})\cong \GL_2(F_{\vv_p})$, and hence local base change is the identity, the local component \(\sigma'_{\vv_p}\) is supercuspidal. However, a non-cuspidal square-integrable automorphic representation cannot have a tempered, and hence {\it a fortiori} supercuspidal, local component, cf.\ \cite{clozel2}, Prop.\ 4.10. Therefore \(\sigma'\) is cuspidal.  
\end{proof}

\begin{thebibliography}{99}
\bibitem[ArtinT]{ArtinTate} E. Artin, J. Tate, \emph{Class Field Theory}, AMS Chelsea Publishing, (2009)
\bibitem[BCGNT]{BCGNT} G. Boxer, F. Calegari, T. Gee, J. Newton, J. Thorne, The Ramanujan and Sato--Tate conjectures for Bianchi modular forms, {\it Forum Math. Pi} \textbf{13} (2025)
\bibitem[Sil94]{Sil94} J. H. Silverman, \emph{Advanced Topics in the Arithmetic of Elliptic Curves}, Graduate Texts in Mathematics, {\bf 151} (1994)
\bibitem[arthur]{arthur-clozel} J. Arthur, L. Clozel, {\it Simple algebras, base change, and the advanced theory of the trace formula}, Ann. Math. Studies, {\bf 120} (Princeton University Press, 1989)
\bibitem[clozel]{clozel2} L. Clozel, On the cohomology of Kottwitz's arithmetic varieties, {\it Duke Math. J.} {\bf 72} (1993) 757--795
\bibitem[hennia]{henniart12} G. Henniart, Induction automorphe globale pour les corps des nombres, {\it Bull. Soc. math. France} {\bf 140} (2012) 1--17
\bibitem[mok]{mok} C. P. Mok, {\it Endoscopic classification of representations of quasi-split unitary groups}, Memoirs of the AMS {\bf 235} (2015).
\bibitem[zou]{zou} R. Chen, J. Zou, Arthur's multiplicity formula for even orthogonal and unitary groups, {\it JEMS} {\bf 27} (2024) 4769--4843
\end{thebibliography}
\end{document}
